\documentclass[11pt]{article}
% Compile with pdfLaTeX from the repository root or from this folder.
\makeatletter
\def\input@path{{../}}
\makeatother
\usepackage{notes}
\setboolean{showanswer}{true}
% Short calculations on the first page need less student writing space.
\renewcommand{\solutionheight}{14mm}

\begin{document}
\lesson[Stewart, Calculus, 9th edition\enspace\textbullet\enspace Section 12.3]
  {1.3}{Dot Product}

\section{A product that measures alignment}
The dot product produces a \textbf{scalar} and connects component calculations with angles.

\begin{definition}[Dot product]
For $\B{u}=\cvec{u_1,u_2,u_3}$ and $\B{v}=\cvec{v_1,v_2,v_3}$,
\[
  \B{u}\cdot\B{v}=u_1v_1+u_2v_2+u_3v_3.
\]
In two dimensions, omit the third term. Multiply corresponding components,
then add the results.
\end{definition}

\begin{example}{Compute $\B{u}\cdot\B{v}$ for
$\B{u}=\cvec{2,-1,3}$ and $\B{v}=\cvec{1,4,-2}$.}
Multiply corresponding components and add:
$\B{u}\cdot\B{v}=2(1)+(-1)(4)+3(-2)=-8$.
\end{example}

\textbf{Properties.} The dot product is commutative and distributes over addition:
\[
 \B{u}\cdot\B{v}=\B{v}\cdot\B{u},\qquad
 \B{u}\cdot(\B{v}+\B{w})=\B{u}\cdot\B{v}+\B{u}\cdot\B{w}.
\]
Scalars factor out: $(c\B{u})\cdot\B{v}=c(\B{u}\cdot\B{v})$.
Also, $\B{u}\cdot\B{u}=\lVert\B{u}\rVert^2$: a vector dotted with itself
gives its squared length.

\section{The angle between vectors}
Place two nonzero vectors tail to tail. Their angle $\theta$ is chosen in
the range $0\leq\theta\leq\pi$. The geometric formula for the dot product is
\[
 \B{u}\cdot\B{v}=\lVert\B{u}\rVert\lVert\B{v}\rVert\cos\theta,
 \qquad
 \cos\theta=\frac{\B{u}\cdot\B{v}}
 {\lVert\B{u}\rVert\lVert\B{v}\rVert}.
\]
A positive dot product means $\theta<90^\circ$; a negative dot product means
$\theta>90^\circ$. At $0$ and $\pi$, the vectors point in the same and opposite
directions, respectively.

\begin{example}{Find the angle between
$\B{u}=\cvec{1,1,0}$ and $\B{v}=\cvec{1,0,1}$.}
We have $\B{u}\cdot\B{v}=1$ and
$\lVert\B{u}\rVert=\lVert\B{v}\rVert=\sqrt{2}$. Thus
$\cos\theta=1/(\sqrt{2}\sqrt{2})=1/2$, giving $\theta=\pi/3=60^\circ$.
\end{example}

\newpage
\renewcommand{\solutionheight}{24mm}
\section{Perpendicular vectors and projections}
\begin{definition}[Orthogonality]
Two vectors are \textbf{orthogonal} when their dot product is zero.
For nonzero vectors, this means the angle between them is $90^\circ$.
The zero vector is orthogonal to every vector, although it has no direction.
\end{definition}

\begin{example}{Are $\B{u}=\cvec{2,-1,3}$ and
$\B{v}=\cvec{1,2,0}$ perpendicular?}
Since $\B{u}\cdot\B{v}=2(1)+(-1)(2)+3(0)=0$, the two nonzero vectors
are perpendicular. No inverse cosine is needed.
\end{example}

\subsection*{The part of a vector in a given direction}
The projection of $\B{a}$ onto a nonzero vector $\B{b}$ is the part of
$\B{a}$ parallel to $\B{b}$. Think of it as the shadow of $\B{a}$ on the
line through $\B{b}$.

\begin{center}
\begin{tikzpicture}[scale=0.85,>=Latex,font=\small]
 \draw[notesrule] (-0.5,0)--(5,0);
 \draw[->,very thick,notesblue] (0,0)--(3,2) node[above]{$\B{a}$};
 \draw[->,thick,notesmuted] (0,-0.9)--(4,-0.9) node[midway,below]{$\B{b}$ direction};
 \draw[->,very thick,notesteal] (0,0)--(3,0)
   node[midway,below]{$\operatorname{proj}_{\B{b}}\B{a}$};
 \draw[dashed,notesmuted] (3,0)--(3,2);
 \draw[notesmuted] (2.8,0)--(2.8,0.2)--(3,0.2);
 \draw[notesink] (0.7,0) arc[start angle=0,end angle=33.69,radius=0.7];
 \node at (0.95,0.3) {$\theta$};
\end{tikzpicture}
\end{center}

\begin{definition}[Scalar and vector projections]
For $\B{b}\ne\B{0}$, the \textbf{scalar projection} of $\B{a}$ onto $\B{b}$ is
\[
 \operatorname{comp}_{\B{b}}\B{a}
 =\frac{\B{a}\cdot\B{b}}{\lVert\B{b}\rVert}.
\]
This is a signed component along $\B{b}$. The \textbf{vector projection} is
\[
 \operatorname{proj}_{\B{b}}\B{a}
 =\frac{\B{a}\cdot\B{b}}{\lVert\B{b}\rVert^2}\B{b}.
\]
\end{definition}

\begin{example}{Project $\B{a}=\cvec{3,1,0}$ onto $\B{b}=\cvec{1,1,0}$.}
Here $\B{a}\cdot\B{b}=4$ and $\lVert\B{b}\rVert^2=2$. Thus
\[
 \operatorname{comp}_{\B{b}}\B{a}=\frac4{\sqrt{2}}=2\sqrt{2},
 \qquad
 \operatorname{proj}_{\B{b}}\B{a}=\frac42\cvec{1,1,0}=\cvec{2,2,0}.
\]
The remaining part is $\B{a}-\operatorname{proj}_{\B{b}}\B{a}=\cvec{1,-1,0}$.
Its dot product with $\B{b}$ is zero, so the remaining part is perpendicular
to the projection direction.
\end{example}

\newpage
\section{Direction angles}
We can describe a nonzero vector's direction using its angles with the
three positive coordinate axes.

\begin{definition}[Direction cosines]
For a nonzero vector $\B{v}=\cvec{v_1,v_2,v_3}$, let $\alpha$, $\beta$, and
$\gamma$ be its angles with the positive $x$-, $y$-, and $z$-axes. Then
\[
 \cos\alpha=\frac{v_1}{\lVert\B{v}\rVert},\qquad
 \cos\beta=\frac{v_2}{\lVert\B{v}\rVert},\qquad
 \cos\gamma=\frac{v_3}{\lVert\B{v}\rVert}.
\]
These \textbf{direction cosines} are the components of the unit vector
$\B{v}/\lVert\B{v}\rVert$.
\end{definition}

For example, $\B{v}\cdot\B{i}=v_1$ and $\lVert\B{i}\rVert=1$, so the angle
formula gives $\cos\alpha=v_1/\lVert\B{v}\rVert$. The other two formulas
follow in the same way. Since the unit direction has length $1$,
\[
 \cos^2\alpha+\cos^2\beta+\cos^2\gamma=1.
\]

\begin{example}{Find the direction angles of $\B{v}=\cvec{1,-1,\sqrt{2}}$.}
Its magnitude is $\lVert\B{v}\rVert=\sqrt{1+1+2}=2$. Therefore,
\[
 \cos\alpha=\frac12,\qquad
 \cos\beta=-\frac12,\qquad
 \cos\gamma=\frac{\sqrt{2}}2.
\]
The direction angles are $\alpha=60^\circ$, $\beta=120^\circ$, and
$\gamma=45^\circ$. The negative $y$-component gives an obtuse angle
with the positive $y$-axis.
\end{example}

\section{Work done by a constant force}
Only the part of a force in the direction of motion contributes to work.
For a constant force $\B{F}$ and a displacement $\B{d}$,
\[
 W=\B{F}\cdot\B{d}
  =\lVert\B{F}\rVert\lVert\B{d}\rVert\cos\theta.
\]
Work is positive when the force has a component along the displacement,
zero when they are perpendicular, and negative when the force has a
component against the displacement.

\begin{example}{A constant force $\B{F}=\cvec{3,4}$ newtons moves an object
from $P=(1,2)$ to $Q=(5,2)$, with coordinates measured in meters.
Find the work done by the force.}
The displacement is $\B{d}=\vec{PQ}=\cvec{4,0}$ meters, so
\[
 W=\B{F}\cdot\B{d}=3(4)+4(0)=12\text{ joules}.
\]
The vertical component of the force contributes no work because the
displacement is horizontal.
\end{example}
\end{document}
